\section{DD3：全局近端存在性}\label{sec:dd3}

DD1 和 DD2 假设所用的近端步存在。本节沿用
Brito、Lara 和 Yuan~\cite[定理~3.1]{BritoLaraYuan2026} 的同伦方法，
先证明引言中的存在引理，再与局部收敛条件结合。

\begin{proof}[引理~\ref{lem:existence} 的证明]
固定任意 $b\in\R^n$ 和 $p_0\in S$，令
\[
 q=b-p_0,\qquad r=\norm q,\qquad \mathcal G(w)=q-cF(p_0+w).
\]
只需证明 $w\in\mathcal G(w)$ 有解。先考察同伦方程
\begin{equation}\label{eq:dd3-homotopy}
 w\in\tau\mathcal G(w),\qquad 0\le\tau\le1.
\end{equation}
当 $\tau=0$ 时 $w=0$。当 $\tau>0$ 时，取相应的 $u\in F(p_0+w)$，记
\[
 y=p_0+w,\quad \mu=\tau c,\quad z=p_0+\tau q=y+\mu u,
 \quad t=d(z)\le r,\quad s=\norm{z-y}.
\]
由锚定条件存在 $p_z\in P_S(z)$，使
\[
 \Phi_\mu(t)\le\ip{z-y}{y-p_z}
 =\ip{z-y}{z-p_z}-s^2\le ts-s^2.
\]
因此实际同伦解满足
\begin{equation}\label{eq:dd3-step-bound}
 s\le\frac{t+\sqrt{t^2-4\Phi_\mu(t)}}2
 \le K_c(r):=\frac{r+\sqrt{r^2+4D_c(r)}}2,
 \qquad \norm w\le r+K_c(r).
\end{equation}
右端有限，且不依赖 $\tau$。

若 $F$ 下半连续且非空闭凸值，则 $\mathcal G$ 也具有这些性质。
Michael 选择定理给出连续选择 $g(w)\in\mathcal G(w)$。
其定义域 $\R^n$ 仿紧，值域是 Banach 空间，且有限维连续映射将有界集映入相对紧集。
方程 $w=\tau g(w)$ 的解均满足上述一致界，故 Schaefer 不动点定理给出
$w=g(w)\in\mathcal G(w)$。
若 $F$ 上半连续且非空紧凸值，则 $\mathcal G$ 同样如此；
上半连续紧值映射在紧集上的值之并为紧集，因此它在有界集上相对紧。
由多值 Leray--Schauder 原理及同一先验界，仍得 $w\in\mathcal G(w)$。
所用三项经典定理见\cite[定理~2.5--2.7]{BritoLaraYuan2026}。
两种情形均得到 $b\in y+cF(y)$，故 $\ran(I+cF)=\R^n$。
\end{proof}

存在性证明不使用正则性函数 $\psi$。条件~\eqref{eq:global-defect}
须在整个参数区间 $(0,c]$ 成立，因为同伦经过其中的每个参数。
若每个 $c>0$ 都满足该条件，则每个正参数下近端步均全局存在。

局部迭代有时还限制输出及所用图值，例如
\begin{equation}\label{eq:dd3-allowed}
 \mathcal J_c(x)=\{y\in J_c(x):y\in V,\ (x-y)/c\in W\}.
\end{equation}
此时须验证这个受限关系本身非空。一个直接的充分条件是：对某个 $p_0\in S$、$r>0$，
\begin{equation}\label{eq:dd3-local-roots}
 V\supseteq B(p_0,R_V),\quad W\supseteq B(0,R_W),\qquad
 r+K_c(r)<R_V,\quad K_c(r)<cR_W.
\end{equation}
事实上，对 $\norm{x-p_0}<r$ 的每个完整近端输出 $y$，
在~\eqref{eq:dd3-step-bound} 中取 $\mu=c,z=x$ 可得
$\norm{x-y}\le K_c(r)$、$\norm{y-p_0}<r+K_c(r)$。
故全部输出都满足限制，即 $\mathcal J_c(x)=J_c(x)\ne\varnothing$。

\begin{theorem}[DD3]\label{thm:dd3}
设引理~\ref{lem:existence} 的条件在参数 $c>0$ 下成立。
取开集 $U$、$R>0$ 及 $U_R=\{x\in U:d(x)\le R\}$。
若使用完整关系 $\mathcal J_c=J_c$，或已验证受限关系
$\mathcal J_c\subseteq J_c$ 在 $U_R$ 上非空，则对初值 $x^0\in U_R$ 有以下结论。
若同一关系满足 DD1 的其余条件及定理~\ref{thm:dd1-trajectory} 的初值条件，
则每条允许的 PPA 轨道均具有有限长度并线性收敛于 $S\cap U$ 中的一点。
若同一关系满足 DD2 及定理~\ref{thm:dd2-series} 或
定理~\ref{thm:dd2-length} 的其余条件与初值条件，
则每条允许轨道均具有相应的距离估计、有限长度及点收敛结论。
\end{theorem}

\begin{proof}
存在引理保证完整近端关系处处非空；受限情形的非空性由附加假设给出。
对同一关系应用所引定理，其有限前缀长度估计保证下一步仍在 $U_R$，
从而可无限续接，并得到相应收敛结论。
\end{proof}

\begin{example}[非二次条件下的全局存在性]\label{ex:dd3-nonquadratic}
设 $Q$ 为平面直角旋转，定义连续算子
\[
 F(y)=-\norm y^{-1/4}y+\norm y^{-1/2}Qy\quad(y\ne0),\qquad F(0)=0.
\]
其零点集为 $S=\{0\}$。记 $r=\norm y$、$t=\norm{y+\mu F(y)}$，则
\begin{equation}\label{eq:dd3-nonquadratic-identities}
 \ip{F(y)}y=-r^{7/4},\qquad
 t^2=(r-\mu r^{3/4})^2+\mu^2r,\qquad
 \norm{F(y)}^2=r^{3/2}+r.
\end{equation}
因此零点附近不存在有限常数 $\gamma$ 使
$\ip{F(y)}y\ge-\gamma\norm y^2$。
另一方面，将 $t^2$ 对实参数 $\mu$ 最小化可得
\[
 t^2\ge\frac{r^2}{1+\sqrt r},\qquad \mu r^{7/4}\le tr^{5/4}\quad(\mu>0).
\]
当 $r\le1$ 时 $r\le\sqrt2\,t$；当 $r\ge1$ 时
$r\le2^{2/3}t^{4/3}$。故所有 $\mu>0$ 共同满足
\[
 \ip{\mu F(y)}y\ge-2^{5/8}t^{9/4}-2^{5/6}t^{8/3}=:\Phi(t).
\]
此下界在有界距离区间内一致有界，存在引理遂给出
$\ran(I+cF)=\R^2$，对每个 $c>0$ 成立。
同时 $d(y)\le r_F(y)^2$。本例说明非二次锚定条件可以保证近端存在性；
收敛仍须另验 DD1 或 DD2 的条件。
\end{example}
